% !PRINTED none
% Part-title (Abschnitt-opening) page: no running head / page number printed on this leaf.

\begin{center}
\gesperrt{Elfter Abschnitt.}

\textbf{\gesperrt{Die Wendepunkte einer Curve dritter Ordnung.}}

\noindent\rule{0.25\linewidth}{0.4pt}

\S.\,85.

\gesperrt{Ternäre Formen und algebraische Curven.}
\end{center}

Um im weiteren Verlauf unserer Darstellung dem Leser
einige von den mannigfaltigen und interessanten Anwendungen
der Algebra auf geometrische Probleme vorführen zu können,
ohne ihn auf andere Hülfsmittel zu verweisen, sollen hier die
nothwendigen geometrischen Grundlagen für diese Anwendungen
in der Kürze abgeleitet werden. Wir beschränken uns dabei
auf die Geometrie der Ebene und lassen auch da Alles bei Seite,
was für unsere Anwendungen nicht direct erforderlich ist\textsuperscript{1)}.

Wir wenden die homogenen Dreieckscoordinaten an, die dadurch erklärt sind, dass die Lage eines Punktes $x$ durch seine
Abstände von den drei Seiten des Coordinatendreiecks, in drei
beliebigen Maasseinheiten gemessen (oder mit drei beliebigen
constanten Factoren multiplicirt), bestimmt sind. Diese drei
Grössen heissen die \gesperrt{Coordinaten} des Punktes $x$, und werden
meist mit $x_1, x_2, x_3$ oder mit $y_1, y_2, y_3$ etc. bezeichnet.

Zwischen diesen drei Coordinaten des Punktes $x$ besteht
eine lineare, nicht homogene Relation, die man erhält, wenn man
den Flächeninhalt des Coordinatendreiecks $(1, 2, 3)$ gleich der
Summe der Inhalte der drei Dreiecke $(x, 2, 3)$, $(x, 3, 1)$, $(x, 1, 2)$
setzt. Eine Coordinate und also auch der Ausdruck für den
Flächeninhalt des entsprechenden der drei letztgenannten Drei-
%CONT>>

\noindent\rule{0.25\linewidth}{0.4pt}

\footnotesize
\noindent\textsuperscript{1)} Ausführlicheres findet man in den Lehrbüchern der analytischen
Geometrie, unter denen das Werk von \gesperrt{George Salmon}, \glqq Higher plane
curves\grqq{}, deutsch von \gesperrt{Fiedler}, hervorzuheben ist.
\normalsize
